
YPF Boxes es el servicio de mantenimiento vehicular integrado a las estaciones de 
servicio. A diferencia de los otros módulos, donde el foco está en cuantificar el 
efecto causal de una intervención específica, este módulo tiene un objetivo 
principalmente de diagnóstico: entender qué factores explican el grado de utilización 
actual de YPF Boxes y en qué contextos operativos alcanzan mayor demanda. Esta 
caracterización es el punto de partida necesario para cualquier decisión futura sobre 
la expansión o reconfiguración del servicio.

\subsection{Objetivo}

Las preguntas centrales que guían el análisis son:

\begin{itemize}
    \item ¿Qué factores explican las diferencias en el grado de utilización entre 
    estaciones?
    \item ¿La presencia de competidores cercanos —lubricentros independientes u 
    otras estaciones con servicios similares- reduce la utilización?
    \item ¿La ubicación de la estación (sobre ruta o en zona urbana) afecta la 
    demanda del servicio?
    \item ¿Existen patrones de estacionalidad o picos de demanda sistemáticos?
    \item ¿Qué características de la estación o del entorno predicen un mayor 
    nivel de utilización?
\end{itemize}

\subsection{Diseño econométrico}

A diferencia de los módulos anteriores, en este caso no se dispone de una 
intervención con variación exógena clara que permita aplicar un diseño 
cuasi-experimental. El análisis se basa en cambio en la variación observada entre 
estaciones y a lo largo del tiempo, estimando un modelo que relaciona el grado de 
utilización de YPF Boxes con un conjunto de factores estructurales, competitivos y 
temporales. El objetivo no es estimar el efecto causal de una política específica, 
sino identificar patrones sistemáticos que permitan caracterizar la demanda y orientar 
decisiones operativas.



\subsubsection*{Ejercicio: Determinantes del grado de utilización}

\textbf{Pregunta:} ¿Qué características de la estación, del entorno competitivo y 
del contexto temporal explican las diferencias en el nivel de uso de los YPF Boxes?

\textbf{Lógica del ejercicio:} Se estima un modelo de regresión donde la variable 
dependiente es el grado de utilización de YPF Boxes en cada estación y período 
—definido como la proporción del tiempo en que los boxes se encuentran ocupados, 
o una métrica equivalente construida a partir de registros operativos—. Entre las 
variables explicativas se incluyen: indicadores de competencia local (cantidad de 
lubricentros y otras estaciones con servicios similares en un radio geográfico 
definido), características de la estación (ubicación sobre ruta o en zona urbana, 
antigüedad, tamaño), variables de precio, y controles temporales para capturar 
estacionalidad y variaciones de corto plazo. La estimación explota variación tanto 
entre estaciones como dentro de cada una a lo largo del tiempo, lo que permite 
separar factores estructurales de fluctuaciones coyunturales.

Los resultados de este ejercicio permiten identificar en qué contextos los YPF Boxes 
operan cerca de su capacidad y en cuáles presentan capacidad ociosa persistente, 
lo que aporta información concreta para decisiones sobre precios, promociones 
o reasignación de recursos dentro de la red.

Los resultados de este bloque persiguen dos objetivos complementarios. En primer lugar, buscan comprender de manera sistemática los patrones de demanda del servicio y el entorno competitivo en el que operan los YPF Boxes. En segundo lugar, procuran sentar las bases para la aplicación futura de las herramientas econométricas desarrolladas en los módulos anteriores, a fin de evaluar potenciales intervenciones orientadas a mejorar el desempeño de esta unidad de negocios.




\subsection{Datos requeridos}

Registros operativos de utilización de YPF Boxes por  estación y período; características de cada estación; información georreferenciada 
sobre competidores. Las fuentes transaccionales requeridas se detallan en la 
sección~\ref{sec:datos}.